% TOP_PROBLEM: {"id": 2829, "title": "Kirby Problem 3.31", "category_id": 7, "category": {"id": 7, "name": "topology", "display_name": "Topology", "description": "Properties preserved under continuous deformations.", "slug": "topology", "order_index": 7, "created_at": "2026-07-31T15:26:25.671Z"}, "status": "open", "problem_number": "KP-3.31", "set_id": 11}
% FIRST_POSTED: 2026-10-01
% ENTRY_KIND: statement_only
% UnsolvedMath ID 2829; source catalogue code KP-3.31
% !TeX program = lualatex
% Definition and sources only; this file is not an LLM solution attempt.
\documentclass[11pt,a4paper]{article}
\usepackage[margin=25mm]{geometry}
\usepackage{fontspec}
\setmainfont{lmroman10-regular.otf}[BoldFont=lmroman10-bold.otf,
 ItalicFont=lmroman10-italic.otf,BoldItalicFont=lmroman10-bolditalic.otf]
\usepackage{amsmath,amssymb,mathtools}
\usepackage{unicode-math}
\setmathfont{latinmodern-math.otf}
\AtBeginDocument{\let\setminus\smallsetminus}
\usepackage{xurl,hyperref}
\hypersetup{colorlinks=true,urlcolor=blue}
\setcounter{secnumdepth}{0}
\setlength{\parindent}{0pt}
\setlength{\parskip}{5pt}
\setlength{\emergencystretch}{3em}
\begin{document}
\section*{Kirby Problem 3.31}\label{2829}
{\small UnsolvedMath ID 2829 · KP-3.31 · Topology · Kirby's Problems in Low-Dimensional Topology}
\subsection{Definitions and mathematical statement}
A nontrivial group $G$ is left-orderable if it admits a strict total order
$<$ satisfying $a<b\Rightarrow ga<gb$ for every $a,b,g\in G$. Equivalently
there is a positive cone $P\subset G$ with
\[
 PP\subseteq P,\qquad G=P\sqcup\{1\}\sqcup P^{-1}.
\]
The displayed union is disjoint. Following the source, exclude the trivial
group from the property; its alternative conventional treatment changes
only that special input case.

Given a finite triangulation of a closed three-manifold $M$, can an algorithm
decide whether $\pi_1(M)$ is left-orderable? It must halt on both positive
and negative instances, with the triangulation providing a finite presentation
for the fundamental group.

Also determine the complexity of certificates of left-orderability or
non-left-orderability. A certificate means finite data whose validity can be
checked from the input; the desired size and checking time should be stated
as functions of its length. Describing an arbitrary infinite positive cone
without a finite effective rule is not a finite certificate. One may also
study the complexity of recognizing positive words for a specified order,
but this is different from deciding whether any left order exists.
\subsection{Short English statement}
Decide left-orderability for closed three-manifold groups and understand how efficiently either answer can be certified.
\subsection{Sources}
\begin{itemize}
\item[S1] ULAM.AI and the UnsolvedMath Contributors, supplied problems.json, record 2829 (KP-3.31), Kirby's Problems in Low-Dimensional Topology. Catalogue identity and source transcription; CC BY 4.0.
\url{https://huggingface.co/datasets/ulamai/UnsolvedMath}.
\item[S2] K3: A New Problem List in Low-Dimensional Topology, Problem 3.31. Expanded formulation of the supplied catalogue statement.
\url{https://math.berkeley.edu/sites/default/files/surv-295-ruberman-watermarked-author-pdf.pdf}.
\end{itemize}
\end{document}
